\documentclass[tikz,border=12pt]{standalone}
\usepackage{braket}
\usepackage{tikz-3dplot}

\tdplotsetmaincoords{70}{115}

\begin{document}
\begin{tikzpicture}[tdplot_main_coords,scale=2.7,>=stealth,font=\small]
  \def\polarAngle{55}
  \def\azimuthAngle{75}

  % Sphere outline and the equatorial reference circle.
  \draw[tdplot_screen_coords,gray!75,line width=0.65pt] (0,0) circle (1);
  \tdplotdrawarc[densely dashed,gray!65,line width=0.55pt]
    {(0,0,0)}{1}{0}{360}{}{}

  % Cartesian axes and the computational basis states.
  \draw[->,gray!85,line width=0.65pt] (0,0,0) -- (1.27,0,0)
    node[below] {$x$};
  \draw[->,gray!85,line width=0.65pt] (0,0,0) -- (0,1.27,0)
    node[right] {$y$};
  \draw[->,gray!85,line width=0.65pt] (0,0,-1.16) -- (0,0,1.25)
    node[above] {$z$};
  \node[left] at (0,0,1) {$\ket{0}$};
  \node[left] at (0,0,-1) {$\ket{1}$};

  % A general pure state and its projection onto the xy plane.
  \tdplotsetcoord{P}{1}{\polarAngle}{\azimuthAngle}
  \draw[densely dashed,gray!75] (P) -- (Pxy);
  \draw[densely dashed,gray!75] (0,0,0) -- (Pxy);
  \draw[->,line width=1.15pt] (0,0,0) -- (P)
    node[above right] {$\ket{\psi}$};

  % Azimuth phi is measured in the xy plane from +x.
  \tdplotdrawarc[->,line width=0.7pt]
    {(0,0,0)}{0.46}{0}{\azimuthAngle}{anchor=north}{$\phi$}

  % Polar angle theta is measured from +z toward the state vector.
  \tdplotsetthetaplanecoords{\azimuthAngle}
  \tdplotdrawarc[tdplot_rotated_coords,->,line width=0.7pt]
    {(0,0,0)}{0.58}{0}{\polarAngle}{anchor=west,xshift=3pt}{$\theta$}
\end{tikzpicture}
\end{document}
